\documentclass[12pt,a4paper]{report}

\usepackage[T1]{fontenc}
\usepackage[utf8]{inputenc}
\usepackage[french]{babel}
\usepackage{lmodern}
\usepackage{geometry}
\usepackage{setspace}
\usepackage{parskip}
\usepackage{hyperref}
\usepackage{enumitem}
\usepackage{listings}
\usepackage{xcolor}

\geometry{margin=2.5cm}
\onehalfspacing

\hypersetup{
  colorlinks=true,
  linkcolor=blue!50!black,
  urlcolor=blue!50!black,
  pdftitle={Rapport synthétique WaveControl},
  pdfauthor={Équipe projet}
}

\lstset{
  basicstyle=\ttfamily\small,
  frame=single,
  breaklines=true,
  columns=fullflexible
}

\begin{document}

\begin{titlepage}
    \centering
    {\scshape\LARGE Rapport synthétique\par}
    \vspace{1cm}
    {\Huge\bfseries WaveControl\par}
    \vspace{0.5cm}
    {\Large Version LaTeX du rapport vierge\par}
    \vspace{1.5cm}

    {\large Date : 4 mars 2026\par}
    \vspace{0.2cm}
    {\large Auteur : Équipe projet\par}

    \vfill
\end{titlepage}

\tableofcontents
\clearpage

\chapter{Choix des outils}

\section{Flutter / Dart}
Le projet a été développé avec Flutter et Dart pour une raison principale : disposer d’une application mobile unique qui fonctionne à la fois sur iOS et Android.

Ce choix permet :
\begin{itemize}[leftmargin=*]
  \item de partager la même base de code,
  \item de réduire le temps de développement,
  \item de simplifier la maintenance.
\end{itemize}

\section{Visual Studio Code}
Visual Studio Code a été utilisé comme environnement de développement car il est léger, pratique pour Flutter, et bien intégré à GitHub Copilot.

\section{Pourquoi GitHub Copilot}
Copilot a été un vrai accélérateur dans ce projet.

Contexte de travail :
\begin{itemize}[leftmargin=*]
  \item le langage Dart n’était pas maîtrisé au départ,
  \item le projet devait être réalisé en environ 3 mois,
  \item il fallait produire rapidement un résultat fonctionnel.
\end{itemize}

Copilot a aidé à :
\begin{itemize}[leftmargin=*]
  \item gagner du temps sur la syntaxe et la structure du code,
  \item avancer malgré la montée en compétence sur un nouveau langage,
  \item rester efficace grâce à une expérience déjà positive avec cet outil.
\end{itemize}

\chapter{Rôle de l’application}

L’application WaveControl a pour rôle de centraliser la gestion du système autour de la base station.

Fonctions principales :
\begin{itemize}[leftmargin=*]
  \item configurer la base station,
  \item ajouter et gérer des configurations,
  \item ajouter et gérer des télécommandes IR,
  \item monitorer l’état du système,
  \item piloter les équipements.
\end{itemize}

En résumé, l’application sert à la fois à configurer, superviser et commander l’installation.

\chapter{Utilisateurs : 3 modes}

L’application est organisée autour de 3 modes d’utilisation :
\begin{itemize}[leftmargin=*]
  \item \textbf{Mode Utilisateur} : accès aux fonctions essentielles de supervision et d’usage quotidien.
  \item \textbf{Mode Technicien} : accès intermédiaire pour la configuration et les réglages opérationnels.
  \item \textbf{Mode Développeur} : accès avancé pour les fonctions techniques et les tests.
\end{itemize}

Cette séparation permet d’adapter l’interface au profil et d’éviter des manipulations sensibles par des utilisateurs non techniques.

\chapter{Mode de communication}

\section{Wi-Fi + Base Station}
Le smartphone communique avec la base station via le réseau Wi-Fi (connectivité IP).

\section{Communication MQTT}
Une fois la connectivité Wi-Fi active, la logique métier de l’application repose sur un service centralisé, \texttt{MQTTService}, implémenté sous forme de singleton et exposé à l’interface via \texttt{ChangeNotifier}. Ce choix permet de concentrer dans une seule couche la connexion au broker, les abonnements, la publication des commandes et la réception des états.

\section{Établissement de session MQTT (dans le code)}
Au démarrage, l’application lance \texttt{MQTTService().connect(...)} puis, lorsque la connexion est confirmée, publie une requête initiale \texttt{publishMessage('home/matter/request', 'test')} afin d’obtenir l’inventaire des équipements. Cette stratégie évite un écran vide au lancement et initialise rapidement la vue de supervision.

La session est construite autour de cinq éléments techniques complémentaires :
\begin{itemize}[leftmargin=*]
  \item \textbf{\texttt{MqttServerClient}} (package \texttt{mqtt\_client}) : composant responsable de l’ouverture de session, des abonnements et des échanges MQTT.
  \item \textbf{\texttt{keepAlivePeriod = 60}} : envoi périodique d’un signal de vie toutes les 60 secondes pour permettre au broker de détecter une perte de session.
  \item \textbf{Authentification utilisateur/mot de passe} : contrôle d’accès au broker pour éviter les connexions et publications non autorisées.
  \item \textbf{Last Will sur \texttt{home/status} avec \texttt{Offline}} : message de présence publié automatiquement par le broker si la session se coupe brutalement.
  \item \textbf{Publication \texttt{Online} après connexion} : signal explicite de disponibilité émis par l’application une fois la session établie.
\end{itemize}

L’association \texttt{Online}/\texttt{Offline} fournit un mécanisme simple et robuste de supervision de présence du client MQTT.

\section{Abonnements et topics suivis}
Après connexion, l’application s’abonne à \texttt{home/\#} (écoute globale) ainsi qu’à des motifs ciblés comme \texttt{home/+/state}, \texttt{home/+/status} et \texttt{home/+/color}. Cette combinaison permet de capter à la fois les réponses d’inventaire et les mises à jour d’état plus fines (ON/OFF, luminosité, couleur), sans limiter l’extensibilité future des équipements.

\section{Publication de commandes depuis l’UI (exemples concrets)}
Les écrans publient leurs actions via des méthodes de service, ce qui évite de disperser la logique protocolaire dans l’interface. Par exemple, \texttt{setBrightness(topic, brightness)} convertit une valeur utilisateur de 0 à 100 vers une valeur 0 à 255 avant publication sur \texttt{.../brightness/set}. De la même manière, \texttt{setRgbColor(topic, r, g, b)} publie sur \texttt{.../color/set} avec un payload RGB, et \texttt{publishMessage(topic, message)} reste disponible pour les commandes génériques.

Dans les parcours métier, plusieurs topics illustrent cette logique : \texttt{home/wristband/request}, \texttt{home/wristband/get\_config}, \texttt{home/wristband/config}, ainsi que \texttt{home/remote/saved}, \texttt{home/remote/new} et \texttt{home/IR/feedback} pour le périmètre infrarouge.

\section{Traitement des messages entrants et mise à jour des états}
La méthode \texttt{\_handleIncomingMessage(topic, payload)} centralise l’interprétation des messages reçus. Deux mécanismes structurent ce traitement :
\begin{enumerate}[leftmargin=*]
  \item pour \texttt{home/matter/response}, le service parse la réponse JSON, met à jour \texttt{deviceStates}, extrait les champs d’état (dont \texttt{brightness} et \texttt{rgb\_color}) et retire les équipements absents de la réponse courante ;
  \item pour les topics incrémentaux (ex. \texttt{home/lum\_1/brightness/state}), le service met à jour uniquement la propriété concernée, puis notifie l’interface via \texttt{notifyListeners()}.
\end{enumerate}

Cette centralisation garantit une cohérence de l’état local et limite les divergences entre écrans.

\section{Canal spécialisé bracelet/IR}
Les flux bracelet et IR sont également conservés dans \texttt{recentMessages} (50 derniers messages) afin de laisser les écrans métiers appliquer leur logique de parsing et de déduplication. Des vues comme \texttt{configuration\_screen.dart}, \texttt{view\_configs\_screen.dart} et \texttt{ir\_device\_detail\_screen.dart} exploitent cette file pour transformer les retours MQTT en comportements UI spécifiques.

\section{Robustesse réseau et reconnexion automatique}
Le service surveille la connectivité via \texttt{connectivity\_plus}. Lorsqu’une coupure est détectée puis résolue, une reconnexion automatique est tentée avec les paramètres mémorisés. Après reconnexion, l’interface est notifiée et une requête \texttt{home/matter/request} est relancée pour resynchroniser l’état global.

En complément, un polling périodique de puissance bracelet est exécuté toutes les 30 secondes sur \texttt{home/wristband/request/power}. La réponse \texttt{home/wristband/feedback/power} est parsée puis exposée via \texttt{batteryStatusNotifier}, ce qui permet un affichage réactif de la batterie et de l’alimentation secteur.

\section{Pourquoi ce design MQTT est pertinent}
Ce design est pertinent car il maintient un point de vérité unique pour les états, conserve une interface réactive (\texttt{ChangeNotifier}/\texttt{ValueNotifier}) et sépare proprement transport MQTT et logique métier d’écran. Cette organisation améliore la robustesse et la maintenabilité, en particulier dans un contexte IoT où les messages sont asynchrones et le réseau potentiellement instable.

\chapter{Organisation du projet (arborescence)}

Présenter l’arborescence du code est pertinent dans un rapport technique, à condition de rester synthétique. L’objectif n’est pas de lister tous les fichiers, mais de montrer la logique d’organisation.

Extrait d’arborescence utile :

\begin{lstlisting}
lib/
  main.dart                 -> point d'entrée de l'application
  screens/                  -> écrans (Home, Settings, Monitoring, Configuration, IR...)
  services/                 -> logique applicative (MQTT, paramètres, connectivité)
  models/                   -> structures de données (états, configurations)
  widgets/                  -> composants UI réutilisables
  theme/                    -> thème visuel (couleurs, styles)
\end{lstlisting}

Pourquoi cette structure est importante :
\begin{itemize}[leftmargin=*]
  \item elle sépare clairement l’interface et la logique métier,
  \item elle facilite la maintenance et l’évolution du projet,
  \item elle aide un nouveau développeur à comprendre rapidement le code.
\end{itemize}

\section{Détail du dossier \texttt{screens/}}
Le dossier \texttt{screens/} regroupe les pages visibles par l’utilisateur et les parcours de navigation.

\begin{itemize}[leftmargin=*]
  \item \texttt{splash\_screen.dart} : écran de démarrage et transition vers l’accueil.
  \item \texttt{home\_screen.dart} : hub principal, état global et accès aux modules selon le mode.
  \item \texttt{settings\_screen.dart} : paramètres généraux (mode, MQTT, préférences).
  \item \texttt{monitoring\_screen.dart} : supervision des équipements et retours d’état.
  \item \texttt{configuration\_screen.dart} : configuration bracelet et associations geste $\rightarrow$ action.
  \item \texttt{view\_configs\_screen.dart} : consultation des configurations existantes.
  \item \texttt{mqtt\_control\_page.dart} : page technique de test MQTT (mode développeur).
  \item \texttt{add\_ir\_screen.dart}, \texttt{IR\_configuration\_screen.dart}, \texttt{periph\_saved\_screen.dart}, \texttt{ir\_device\_detail\_screen.dart} : gestion des télécommandes IR (ajout, liste, détail, maintenance).
\end{itemize}

En résumé, \texttt{screens/} porte l’expérience utilisateur et la logique de navigation.

\section{Détail du dossier \texttt{services/}}
Le dossier \texttt{services/} centralise la logique transverse et la communication.

\begin{itemize}[leftmargin=*]
  \item \texttt{mqtt\_service.dart} : connexion au broker MQTT, abonnement aux topics, publication des commandes, réception des messages, exposition des états utiles à l’UI.
  \item \texttt{app\_settings.dart} : gestion des préférences persistantes (mode utilisateur, paramètres de connexion, thème, langue, etc.).
\end{itemize}

Ce dossier est le cœur technique de l’application : les écrans consomment ces services au lieu de gérer directement le réseau et le stockage.

\section{Détail du dossier \texttt{models/}}
Le dossier \texttt{models/} contient les structures de données métier manipulées dans les écrans et services.

\begin{itemize}[leftmargin=*]
  \item \texttt{device\_state.dart} : représentation d’un équipement (identifiant, état, propriétés d’affichage, métadonnées).
  \item \texttt{wristband\_config.dart} : structures liées aux configurations du bracelet (gestes, actions, statuts).
  \item \texttt{config\_item.dart} : modèle d’élément de configuration utilisé dans les flux de paramétrage.
  \item \texttt{command\_history.dart} : modèle d’historique des commandes/messages pour le suivi et le diagnostic.
\end{itemize}

Le rôle des modèles est de standardiser les données échangées entre UI et services, afin de limiter les erreurs et faciliter l’évolution du code.

\chapter{Conclusion}

Le choix Flutter/Dart, combiné à Visual Studio Code et Copilot, a permis de livrer une application mobile multiplateforme dans un délai court.

WaveControl répond au besoin métier principal : configurer la base station, gérer les configurations et télécommandes IR, monitorer le système et piloter les équipements, avec une architecture d’usage claire basée sur 3 modes et une communication Wi-Fi/MQTT efficace.

\end{document}
